% \section*{\Huge Abstract}

% Trajectory optimization offers mature tools for motion planning in high-dimensional spaces under dynamic constraints. However, in obstacle-cluttered environments, the non-convexity of the free-space typically forces roboticists to rely on sampling-based planners, which struggle with high dimensions and differential constraints. Here, we demonstrate that convex optimization can reliably solve these problems through the convex decomposition of the free space.
% Specifically, we decompose the collision-free configuration space into finite convex regions, organizing them into a Graph of Convex Sets (GCS). By combining Bézier curves with this graph structure, we formulate the motion planning problem as a compact mixed-integer convex program that naturally handles collision avoidance, velocity, and duration constraints.
% We validate GCS across a spectrum of environments with increasing complexity, ranging from simple 2D obstacles and cluttered mazes to dynamical quadrotors and high-dimensional 7-DoF manipulators. Numerical experiments demonstrate that GCS consistently outperforms widely-used sampling-based planners (e.g., PRM). 
% Specifically, in high-dimensional tasks, GCS reduces online query time by up to an order of magnitude compared to standard PRM, while finding globally optimal trajectories that are significantly shorter (up to 40\% reduction in length) than those from post-processed sampling-based planners~\cite{motionplanning2022}. Ultimately, this work brings the reliability and speed of convex optimization to the traditionally challenging domain of nonconvex motion planning.


\section*{\Huge Tóm tắt}

Hoạch định chuyển động cho rô-bốt trong môi trường chứa vật cản là một bài toán khó do không gian tự do thường có cấu trúc phi lồi, trong khi quỹ đạo cần đồng thời thỏa mãn các yêu cầu về an toàn hình học, độ trơn và khả thi động lực học.
Báo cáo này nghiên cứu một khung hoạch định dựa trên phân hoạch không gian tự do thành các vùng lồi và tổ chức chúng dưới dạng Đồ thị các Tập Lồi (Graph of Convex Sets~--- GCS).
Trên cơ sở đó, chúng tôi phát triển và đánh giá hai hướng tiếp cận bổ trợ cho nhau, mỗi hướng nhắm tới một miền bài toán khác biệt.

Hướng thứ nhất, GCS-Bézier, tham số hóa quỹ đạo trong mỗi vùng lồi bằng đường cong Bézier bậc~6 với tính liên tục $C^2$.
Nhờ tính chất convex hull của Bézier và cấu trúc luồng nhị phân trên đồ thị, bài toán chọn đường đi rời rạc và tối ưu hóa quỹ đạo hình học được mô hình hóa thành MISOCP với relaxation lồi chặt.
Thực nghiệm trên 10 môi trường maze 25$\times$25 ngẫu nhiên (hơn 100 vùng lồi mỗi trường hợp) cho thấy thời gian giải trung bình dưới 2~giây, và relaxation gap bằng~0 đạt được tự nhiên trên toàn bộ 10 instance mà không cần chiến lược làm tròn bổ sung.
Về ảnh hưởng của phương pháp phân hoạch, thuật toán Approximate Convex Decomposition (ACD) giảm thời gian giải \textbf{64~lần} so với Visibility Clique Cover (VCC) trên cùng môi trường, với chi phí đường đi chênh lệch dưới 1\%, xác lập ACD là lựa chọn thực tiễn tốt nhất trong số các phương pháp được khảo sát.

Hướng thứ hai, GCS-DMS, mở rộng khung GCS sang bài toán Điều khiển Tối ưu (Optimal Control Problem~--- OCP) kết hợp với kỹ thuật direct Multiple Shooting (DMS).
Thay vì chỉ tối ưu hình dạng quỹ đạo, GCS-DMS mô hình hóa trực tiếp phương trình động lực học phi tuyến $\dot x = f(x,u)$, tín hiệu điều khiển $u(\cdot)$, và thời lượng di chuyển trong từng vùng; ràng buộc an toàn được xử lý bằng hard constraint lưới hữu hạn hoặc hàm cản logarit (log-barrier).
Mô hình thuộc lớp MIOCP ở cấp liên tục và MINLP phi lồi sau rời rạc hóa.
Các thực nghiệm với xe unicycle nonholonomic trên maze 15$\times$15 cho thấy ràng buộc defect động lực học đạt $2.47\times10^{-11}$ và sai số ghép nối giữa các vùng đạt $1.05\times10^{-10}$, chứng minh tính khả thi số học của quỹ đạo thu được.
Về chi phí tính toán, cơ chế log-barrier yêu cầu trung bình khoảng 268~giây trong khi ràng buộc lưới hữu hạn điểm chỉ cần khoảng 72~giây; đổi lại, log-barrier đảm bảo quỹ đạo nằm sâu bên trong vùng an toàn thay vì chỉ kiểm tra điều kiện tại hữu hạn điểm rời rạc.

Nhìn chung, GCS-Bézier cho thấy ưu thế rõ rệt về tốc độ và bảo đảm tối ưu lồi trong các bài toán thiên về ràng buộc hình học, trong khi GCS-DMS cung cấp mô hình giàu ý nghĩa vật lý hơn khi tính khả thi động lực học và ràng buộc điều khiển của hệ thực là yêu cầu trung tâm.
Hai hướng tiếp cận cùng minh họa rằng tối ưu hóa lồi và hỗn hợp nguyên có thể giải quyết đáng tin cậy lớp bài toán hoạch định chuyển động phi lồi vốn trước đây phụ thuộc chủ yếu vào các phương pháp lấy mẫu.
